\section{Resumen del capítulo}

\begin{tcolorbox}[resumen,title=Sistemas de numeración]
Un sistema de numeración permite representar cantidades mediante símbolos y reglas.

En el sistema decimal utilizamos diez cifras:
\[
0,1,2,3,4,5,6,7,8,9.
\]
El valor de una cifra depende de la posición que ocupa. Por ejemplo,
\[
4387=4000+300+80+7.
\]
Por esta razón decimos que el sistema decimal es \emph{posicional}.
\end{tcolorbox}

\begin{tcolorbox}[resumen,title=Números naturales y enteros]
Los números naturales se representan mediante
\[
\mathbb{N}=\{0,1,2,3,\ldots\}.
\]
Los números enteros amplían este conjunto incorporando los números negativos:
\[
\mathbb{Z}
=
\{\ldots,-3,-2,-1,0,1,2,3,\ldots\}.
\]
Por lo tanto,
\[
\mathbb{N}\subseteq\mathbb{Z}.
\]
Todo número entero \(a\) posee un opuesto, representado por \(-a\), que cumple
\[
a+(-a)=0.
\]
\end{tcolorbox}

\begin{tcolorbox}[resumen,title=Adición]
La adición es una operación binaria:
\[
+\colon\mathbb{Z}\times\mathbb{Z}\to\mathbb{Z}.
\]
Dados dos números \(a\) y \(b\), su suma se representa mediante
\[
a+b.
\]
Los números que participan de la suma reciben el nombre de \emph{sumandos}.
Una expresión como
\[
302+22
\]
ya representa un número. Evaluarla consiste en encontrar otra representación equivalente de ese mismo número:
\[
302+22=324.
\]
Cuando aparecen más de dos sumandos, la adición continúa realizándose de a dos. La propiedad asociativa permite omitir los paréntesis cuando solo intervienen adiciones.
\end{tcolorbox}

\begin{tcolorbox}[resumen,title=Propiedades de la adición]
Para cualesquiera enteros \(a\), \(b\) y \(c\), la adición cumple:
  
\textit{Conmutatividad}
\[
a+b=b+a.
\]
\textit{Asociatividad}
\[
(a+b)+c=a+(b+c).
\]
\textit{Elemento neutro}
\[
a+0=a.
\]
\textit{Existencia del opuesto}
\[
a+(-a)=0.
\]

Estas propiedades permiten reorganizar y simplificar sumas sin modificar el número que representan.
\end{tcolorbox}

\begin{tcolorbox}[resumen,title=Sustracción]
La sustracción puede definirse mediante la adición y el opuesto:
\[
a-b=a+(-b).
\]
Por ejemplo,
\(
8-5=8+(-5).
\)

El símbolo \(-\) puede tener dos usos diferentes. En \(-5\) indica el opuesto de \(5\), mientras que en \(8-5\) indica una sustracción.

Una vez reescritas las sustracciones como sumas de opuestos, podemos trabajar utilizando las propiedades de la adición. Por ejemplo,
\[
7-4+2
=
7+(-4)+2.
\]
En esta última expresión los sumandos son \(7\), \(-4\) y \(2\).
\end{tcolorbox}

\begin{tcolorbox}[resumen,title=Igualdades y ecuaciones]
Una igualdad indica que las expresiones situadas a ambos lados de \(=\) representan el mismo número.

Por ejemplo,
\[
3+4=7
\]
y
\[
7=3+4
\]
expresan la misma igualdad.

Una ecuación es una igualdad que contiene uno o más valores desconocidos. Por ejemplo,
\[
x+3=8.
\]
Resolver una ecuación consiste en encontrar los valores que hacen verdadera la igualdad.

Para conservar una igualdad debemos realizar la misma transformación en ambos miembros. Por ejemplo,
\[
x+3=8
\]
puede transformarse sumando \(-3\) en ambos lados:
\[
x+3+(-3)=8+(-3),
\]
de donde
\[
x=5.
\]
La idea de ``pasar un término al otro lado cambiando el signo'' es una abreviatura de estas transformaciones y no una regla independiente.
\end{tcolorbox}
